\documentclass[manuscript]{aomart}

\usepackage{amsmath,amssymb,amsthm,mathtools}
\usepackage{hyperref}
\usepackage{booktabs}

\hypersetup{
    colorlinks=true,
    linkcolor=blue,
    urlcolor=blue,
    citecolor=blue
}

% Macros
\newcommand{\Z}{\mathbb{Z}}
\newcommand{\Q}{\mathbb{Q}}
\newcommand{\R}{\mathbb{R}}
\newcommand{\N}{\mathbb{N}}

% Theorem environments (aomart provides amsthm)
\theoremstyle{plain}
\newtheorem{theoremA}{Theorem}
\renewcommand{\thetheoremA}{\Alph{theoremA}}
\newtheorem{theorem}{Theorem}[section]
\newtheorem{proposition}[theorem]{Proposition}
\newtheorem{lemma}[theorem]{Lemma}
\newtheorem{corollary}[theorem]{Corollary}

\theoremstyle{definition}
\newtheorem{definition}[theorem]{Definition}
\newtheorem{example}[theorem]{Example}

\theoremstyle{remark}
\newtheorem{remark}[theorem]{Remark}

\numberwithin{equation}{section}

%------------------------------------------------------------------
\title{Diophantine Confinement of the Branch Locus\\in Syracuse Dynamics}

\author{John A.\ Janik}
\address{john.janik@gmail.com}

\subjclass[2020]{Primary 11B83; Secondary 37A45, 37B10}
\date{\today}

\keywords{Collatz conjecture, correction ratio, golden mean shift,
  subshift of finite type, cycle elimination, Steiner cycle equation,
  irrationality measure, Lean~4 formalization}

\begin{document}

\begin{abstract}
We prove that if the transverse walk associated to a Collatz
trajectory has eventually linear drift, then the trajectory is
eventually bounded and periodic.  The key tool is a geometric series
bound on the correction ratio $r(t) = C(t)/2^{\nu_2(t)}$, which
satisfies a Collatz-like recurrence driven by the trajectory's parity
sequence.  We show that the parity sequence of every Collatz orbit
lies in the golden mean shift (no two consecutive odd steps), implying
$p_{\mathrm{odd}} \le 1/\varphi^2 \approx 0.382$, within $0.005$ of
the equilibrium threshold $p_{\mathrm{eq}} \approx 0.387$.  Cycle
elimination is proved for all cycles with $\Delta_3 \le 79$ odd steps
per period using Steiner's correction bounds, Rhin's effective
irrationality measure for $\log_2 3$~\cite{Rhin1987}, and Hercher's
computational verification for $m$-cycles with
$m \le 91$~\cite{Hercher2024}.  Combined with a conditional
hypothesis for longer cycles ($\Delta_3 \ge 80$), this yields a
conditional proof of the Collatz conjecture: linear walk drift implies
$\mathrm{collatzReaches}(n)$.  We also prove the converse:
reaching~$1$ implies the K-bound, so the sole critical-path axiom is
formally equivalent to the conjecture itself.  A Syracuse framework
bridge connects Tao's odd-step formulation to the walk/drift
infrastructure.  The argument is formalized in approximately
7\,100~lines of Lean~4 across 30~files with five external axioms
(Baker~1966, Rhin~1987, Hercher~2023, Weyl~1916, and one open
mixing axiom on the $(2,3)$-solenoid), all identified and categorized.
\end{abstract}

\maketitle

%==================================================================
\section{Introduction}\label{sec:intro}
%==================================================================

\subsection{The conjecture}

\begin{definition}\label{def:collatz}
The \emph{Collatz map} $T\colon \N \to \N$ is defined by
\[
  T(n) =
  \begin{cases}
    n/2 & \text{if } n \equiv 0 \pmod{2},\\
    3n+1 & \text{if } n \equiv 1 \pmod{2}.
  \end{cases}
\]
The \emph{Collatz conjecture} asserts that for every $n \ge 1$,
there exists $k \ge 0$ such that $T^k(n) = 1$.
\end{definition}

The conjecture has been verified computationally for all
$n \le 10^{20}$ \cite{Oliveira2015} and remains one of the most
resistant open problems in elementary number theory.  Surveys of the
extensive literature can be found in Lagarias~\cite{Lagarias1985,
Lagarias2010}.

\subsection{Prior work}

Terras~\cite{Terras1976} introduced the stopping time and proved that
almost all integers (in the sense of natural density) have a finite
stopping time.  Lagarias~\cite{Lagarias1985} formalized the
probabilistic heuristic that treats the Collatz iteration as a random
walk with drift $\log_2(4/3)$ per odd step.  Wirsching~\cite{Wirsching1998}
developed the ergodic-theoretic perspective, studying transfer
operators associated to the dynamics.

Steiner~\cite{Steiner1977} and Simons--de Weger~\cite{SimonsDeWeger2005}
used Baker-type bounds on linear forms in logarithms to establish
lower bounds on the period of any nontrivial Collatz cycle.
Eliahou~\cite{Eliahou1993} showed that no nontrivial cycle has period
less than $17\,087\,915$.  Conway~\cite{Conway1972} proved that
Collatz-type problems can be algorithmically undecidable in general,
though this does not apply to the specific $3n+1$ function.

The strongest density result is due to Tao~\cite{Tao2022}, who proved
that almost all orbits (in a logarithmic density sense) attain almost
bounded values.  His approach uses entropy-based methods on the
Syracuse random variables and does not yield the full conjecture.

\subsection{Our approach}

The present paper introduces a different mechanism.  The argument
proceeds in three steps:
\begin{enumerate}
  \item The parity sequence of every Collatz trajectory lies in the
    \emph{golden mean shift} (the subshift of finite type forbidding
    the bigram ``11'').  This constrains the odd-step proportion to
    $p_{\mathrm{odd}} \le 1/\varphi^2 \approx 0.382$, just below the
    equilibrium threshold $p_{\mathrm{eq}} = 1/(1 + \log_2 3)
    \approx 0.387$.
  \item The \emph{correction ratio} $r(t) = C(t)/2^{\nu_2(t)}$,
    which captures the cumulative effect of the ``$+1$'' in $3n+1$,
    satisfies a geometric series bound when the transverse walk has
    linear drift.  This forces the trajectory to be eventually
    bounded.
  \item Cycle elimination---proved directly for $\Delta_3 = 1$,
    and via Steiner correction bounds combined with Hercher's
    computational verification~\cite{Hercher2024} for
    $\Delta_3 \le 79$---shows the only positive cycle is
    $\{1, 2, 4\}$ (conditionally on extending the Hercher bound
    beyond $m = 91$ for $\Delta_3 \ge 80$).
\end{enumerate}

\subsection{Main results}

We state two main theorems.  Theorem~A is unconditional given the
drift hypothesis; Theorem~B is a conditional proof of the full
conjecture.

\begin{theoremA}[Walk drift implies reaches~1]\label{thm:A}
Let $n \ge 1$.  If the transverse walk $w(t)$ has eventually linear
drift with rate $\varepsilon > 0$, and the only positive Collatz
cycle is $\{1, 2, 4\}$, then $T^k(n) = 1$ for some $k \ge 0$.
\end{theoremA}

\begin{theoremA}[Conditional Collatz conjecture]\label{thm:B}
Assume:
\begin{enumerate}
  \item[\textup{(H1)}] For every $n \ge 1$, there exist
    $\varepsilon > 0$, $T_0$, and $K$ such that
    $\nu_3(t)/t \le p_{\mathrm{eq}} - \varepsilon$ and
    $3\nu_3(t) \le t + K$ for all $t \ge T_0$.
  \item[\textup{(H2)}] Hercher's theorem: no nontrivial $m$-cycle
    exists for $m \le 91$~\cite{Hercher2024}.
  \item[\textup{(H3)}] No nontrivial cycle exists with
    $\Delta_3 \ge 80$ odd steps per period.
\end{enumerate}
Then for every $n \ge 1$, there exists $k \ge 0$ such that
$T^k(n) = 1$.
\end{theoremA}

Hypothesis~(H1) is a quantitative version of the statement that the
transverse walk has positive linear drift.  The golden mean shift
constraint makes this plausible: the gap between
$1/\varphi^2 \approx 0.382$ and $p_{\mathrm{eq}} \approx 0.387$ is
only $0.005$.  What remains is to promote this topological bound to a
pointwise ergodic statement for individual trajectories---essentially,
an ergodic theorem for the Collatz transfer operator, or a spectral
gap argument.  Hypothesis~(H2) is computationally verified
mathematics~\cite{Hercher2024}; we treat it as an external axiom in
the formalization.  Hypothesis~(H3) is the frontier: it requires
extending Hercher's verification beyond $m = 91$, or developing
analytical bounds (via Baker's theorem~\cite{Baker1975} and Rhin's
irrationality measure~\cite{Rhin1987}) sufficient to cover all
$\Delta_3 \ge 80$.

\subsection{Outline}

Section~\ref{sec:winding} introduces the winding numbers, transverse
walk, and the multiplicative identity that relates trajectory values
to their parity counts.  Section~\ref{sec:golden} establishes that
every Collatz parity sequence lies in the golden mean shift.
Section~\ref{sec:correction} contains the core result: the correction
ratio bound.  Section~\ref{sec:cycles} treats cycle elimination,
including the Steiner correction bounds and the Hercher computational
verification.  Section~\ref{sec:main} assembles the main theorems.
Section~\ref{sec:lean} describes the Lean~4 formalization.
Section~\ref{sec:computation} summarizes computational evidence.

%==================================================================
\section{Winding numbers and the multiplicative identity}
\label{sec:winding}
%==================================================================

\begin{definition}\label{def:parity-counts}
For $n \ge 1$, define the \emph{Collatz sequence} by $a_0 = n$ and
$a_{t+1} = T(a_t)$.  Let $\nu_2(t)$ and $\nu_3(t)$ count the even
and odd steps among $0, 1, \ldots, t-1$, so that
$\nu_2(t) + \nu_3(t) = t$.  The \emph{transverse walk} is
\begin{equation}\label{eq:walk}
  w(t) = \nu_2(t) - \log_2 3 \cdot \nu_3(t).
\end{equation}
\end{definition}

The quantity $w(t)$ measures the excess of even steps over the
equilibrium rate $\log_2 3$ per odd step.  If $w(t) \to +\infty$,
then even steps dominate sufficiently to force the trajectory toward
small values.

\begin{proposition}[Multiplicative identity]\label{prop:identity}
For all $n \ge 1$ and $t \ge 0$,
\begin{equation}\label{eq:identity}
  a_t \cdot 2^{\nu_2(t)} = n \cdot 3^{\nu_3(t)} + C(t),
\end{equation}
where the \emph{correction term} $C(t)$ satisfies $C(0) = 0$ and
\begin{equation}\label{eq:C-recurrence}
  C(t+1) =
  \begin{cases}
    C(t) & \text{if step } t \text{ is even,}\\
    3\,C(t) + 2^{\nu_2(t)} & \text{if step } t \text{ is odd.}
  \end{cases}
\end{equation}
\end{proposition}

\begin{proof}
By induction on $t$.  The base case $t = 0$ gives
$a_0 \cdot 2^0 = n \cdot 3^0 + 0$.

If step $t$ is even, then $a_{t+1} = a_t/2$,
$\nu_2(t{+}1) = \nu_2(t) + 1$, $\nu_3(t{+}1) = \nu_3(t)$, and
$C(t{+}1) = C(t)$.  The identity at $t{+}1$ becomes
\[
  \frac{a_t}{2} \cdot 2^{\nu_2(t)+1}
  = a_t \cdot 2^{\nu_2(t)}
  = n \cdot 3^{\nu_3(t)} + C(t).
\]

If step $t$ is odd, then $a_{t+1} = 3a_t + 1$,
$\nu_2(t{+}1) = \nu_2(t)$, $\nu_3(t{+}1) = \nu_3(t) + 1$, and
$C(t{+}1) = 3\,C(t) + 2^{\nu_2(t)}$.  The identity at $t{+}1$ becomes
\begin{align*}
  (3a_t + 1) \cdot 2^{\nu_2(t)}
  &= 3 \cdot a_t \cdot 2^{\nu_2(t)} + 2^{\nu_2(t)} \\
  &= 3\bigl(n \cdot 3^{\nu_3(t)} + C(t)\bigr) + 2^{\nu_2(t)} \\
  &= n \cdot 3^{\nu_3(t)+1} + 3\,C(t) + 2^{\nu_2(t)}. \qedhere
\end{align*}
\end{proof}

\begin{definition}\label{def:correction-ratio}
The \emph{correction ratio} is
\[
  r(t) = \frac{C(t)}{2^{\nu_2(t)}}.
\]
\end{definition}

\begin{lemma}[Correction ratio recurrence]\label{lem:r-recurrence}
The correction ratio satisfies:
\begin{itemize}
  \item Even step: $r(t{+}1) = r(t)/2$.
  \item Odd step: $r(t{+}1) = 3\,r(t) + 1$.
\end{itemize}
\end{lemma}

\begin{proof}
At an even step, $\nu_2(t{+}1) = \nu_2(t) + 1$ and
$C(t{+}1) = C(t)$, so
$r(t{+}1) = C(t)/2^{\nu_2(t)+1} = r(t)/2$.

At an odd step, $\nu_2(t{+}1) = \nu_2(t)$ and
$C(t{+}1) = 3\,C(t) + 2^{\nu_2(t)}$, so
$r(t{+}1) = (3\,C(t) + 2^{\nu_2(t)})/2^{\nu_2(t)} = 3\,r(t) + 1$.
\end{proof}

\begin{proposition}\label{prop:a-decomposition}
For all $n \ge 1$ and $t \ge 0$,
\begin{equation}\label{eq:a-decomposition}
  a_t = n \cdot 2^{-w(t)} + r(t).
\end{equation}
In particular, if $w(t) \to +\infty$, then $a_t \approx r(t)$ for
large~$t$.
\end{proposition}

\begin{proof}
Dividing~\eqref{eq:identity} by $2^{\nu_2(t)}$ gives
$a_t = n \cdot 3^{\nu_3(t)} / 2^{\nu_2(t)} + r(t)$.  Since
$3^{\nu_3} / 2^{\nu_2} = 2^{\nu_3 \log_2 3 - \nu_2} = 2^{-w(t)}$,
the result follows.
\end{proof}

%==================================================================
\section{The golden mean shift}\label{sec:golden}
%==================================================================

\begin{proposition}\label{prop:odd-gives-even}
If $n$ is odd, then $3n + 1$ is even.
\end{proposition}

\begin{proof}
If $n$ is odd, then $3n$ is odd, so $3n + 1$ is even.
\end{proof}

\begin{corollary}\label{cor:no-11}
In any Collatz trajectory, no two consecutive steps can both be odd
steps.
\end{corollary}

\begin{proof}
An odd step is applied when $a_t$ is odd, producing
$a_{t+1} = 3a_t + 1$, which is even by
Proposition~\ref{prop:odd-gives-even}.  Therefore step $t{+}1$ is an
even step.
\end{proof}

\begin{definition}\label{def:parity-seq}
The \emph{parity sequence} of the Collatz orbit of $n$ is the
sequence $\sigma = (\sigma_0, \sigma_1, \sigma_2, \ldots) \in
\{0,1\}^{\N}$ where $\sigma_t = a_t \bmod 2$.  The \emph{golden mean
shift} $X_\varphi \subset \{0,1\}^{\N}$ is the subshift of finite
type defined by the forbidden bigram $11$:
\[
  X_\varphi = \bigl\{ \sigma \in \{0,1\}^{\N} :
  \sigma_t \cdot \sigma_{t+1} = 0 \text{ for all } t \ge 0 \bigr\}.
\]
\end{definition}

\begin{theorem}\label{thm:collatz-sft}
The parity sequence of every Collatz orbit lies in the golden mean
shift $X_\varphi$.
\end{theorem}

\begin{proof}
This is a restatement of Corollary~\ref{cor:no-11}: the bigram
$\sigma_t = 1$, $\sigma_{t+1} = 1$ would require two consecutive odd
values, which is impossible.
\end{proof}

\begin{corollary}\label{cor:p-odd-bound}
For any Collatz trajectory of length $t$, the odd-step proportion
satisfies $\nu_3(t)/t \le \lceil t/2 \rceil / t$.  In particular,
$\limsup_{t \to \infty} \nu_3(t)/t \le 1/2$.

More precisely, the topological entropy of the golden mean shift is
$\log \varphi$, where $\varphi = (1+\sqrt{5})/2$ is the golden
ratio.  The measure of maximal entropy assigns weight
$1/\varphi^2 = (3-\sqrt{5})/2 \approx 0.382$ to the symbol~$1$.
Therefore, under the maximal-entropy measure,
$p_{\mathrm{odd}} = 1/\varphi^2$.
\end{corollary}

\begin{proposition}\label{prop:log23-lt-phi}
$\log_2 3 < \varphi$.
\end{proposition}

\begin{proof}
We show $\log_2 3 < 8/5 < \varphi$.  For the first inequality,
$3^5 = 243 < 256 = 2^8$, so $5 \log_2 3 < 8$, giving
$\log_2 3 < 8/5$.  For the second, $(8/5)^2 = 64/25 < 5 = \varphi +
1/\varphi$; more precisely, $\varphi^2 = \varphi + 1 > 8/5 + 1 =
13/5$, and $(8/5)^2 = 64/25 < 65/25 = 13/5$.
\end{proof}

\begin{remark}\label{rem:near-miss}
The equilibrium threshold for the walk~\eqref{eq:walk} is
\[
  p_{\mathrm{eq}} = \frac{1}{1 + \log_2 3} \approx 0.38685.
\]
At this odd-step proportion, the expected walk increment is zero.  The
golden mean shift bound gives
$p_{\mathrm{odd}} \le 1/\varphi^2 \approx 0.38197$.  The gap is
$p_{\mathrm{eq}} - 1/\varphi^2 \approx 0.005$, tantalizingly small.
This means the SFT constraint \emph{nearly} suffices to prove
positive drift.  The remaining challenge is to promote the topological
bound $1/\varphi^2$ to a pointwise ergodic statement for individual
trajectories.  Proposition~\ref{prop:log23-lt-phi} ensures that the
inequality goes the right way: since $\log_2 3 < \varphi$, the
reciprocal satisfies $1/(1+\log_2 3) > 1/(1+\varphi) = 1/\varphi^2$,
confirming that the golden mean shift maximum $1/\varphi^2$ lies
strictly below $p_{\mathrm{eq}}$.
\end{remark}

%==================================================================
\section{The correction ratio bound}\label{sec:correction}
%==================================================================

This section contains the core new result.

\begin{definition}\label{def:linear-drift}
We say the walk has \emph{eventually linear drift} with rate
$\varepsilon > 0$ if there exists $T_0$ such that
$w(t) \ge \varepsilon\,t$ for all $t \ge T_0$.
\end{definition}

\subsection{Unrolling the correction recurrence}

Let the odd steps occur at times $\tau_1 < \tau_2 < \tau_3 < \cdots$.
Write $U_j = \nu_2(\tau_j)$ for the cumulative even-step count at the
$j$-th odd step.  Since even steps leave $C$ unchanged and each odd
step sends $C \mapsto 3C + 2^{U_j}$, unrolling gives:

\begin{lemma}[Correction sum]\label{lem:C-sum}
After the $k$-th odd step,
\begin{equation}\label{eq:C-sum}
  C(\tau_k + 1) = \sum_{j=1}^{k} 3^{k-j} \cdot 2^{U_j}.
\end{equation}
\end{lemma}

\begin{proof}
By induction.  At $k = 1$: $C(\tau_1 + 1) = 3 \cdot 0 + 2^{U_1} =
2^{U_1}$.  Assuming the formula at $k$, the $(k{+}1)$-th odd step
gives $C(\tau_{k+1} + 1) = 3 \cdot C(\tau_k + 1) + 2^{U_{k+1}}$,
which expands to
$\sum_{j=1}^{k} 3^{k+1-j} \cdot 2^{U_j} + 2^{U_{k+1}}
= \sum_{j=1}^{k+1} 3^{k+1-j} \cdot 2^{U_j}$.
\end{proof}

At times just after the $k$-th odd step (where $\nu_2 = U_k$ and
$\nu_3 = k$), the correction ratio is
\begin{equation}\label{eq:r-sum}
  r = \frac{C(\tau_k + 1)}{2^{U_k}}
  = \sum_{m=0}^{k-1} 3^m \cdot 2^{-(U_k - U_{k-m})}
  = \sum_{m=0}^{k-1} 3^m \cdot 2^{-D_m},
\end{equation}
where $D_m = U_k - U_{k-m}$ is the number of even steps between the
$(k{-}m)$-th and $k$-th odd steps.

\subsection{The walk-drift gap bound}

\begin{lemma}[Walk-drift gap bound]\label{lem:D-bound}
If the walk has eventually linear drift with rate $\varepsilon > 0$,
then for all sufficiently large $k$ and all $0 \le m \le k$ with
$\tau_{k-m} \ge T_0$,
\begin{equation}\label{eq:D-bound}
  D_m \ge m \cdot \alpha,
  \qquad\text{where}\quad
  \alpha = \frac{\log_2 3 + \varepsilon}{1 - \varepsilon} > \log_2 3.
\end{equation}
\end{lemma}

\begin{proof}
Apply the drift condition at times $\tau_{k-m}$ and $\tau_k$:
\begin{align*}
  w(\tau_k) &= U_k - \log_2 3 \cdot k \ge \varepsilon\,\tau_k,\\
  w(\tau_{k-m}) &= U_{k-m} - \log_2 3 \cdot (k{-}m) \ge \varepsilon\,\tau_{k-m}.
\end{align*}
Subtracting gives
$D_m - \log_2 3 \cdot m \ge \varepsilon\,(\tau_k - \tau_{k-m})$.
Between the $(k{-}m)$-th and $k$-th odd steps, the trajectory takes
$m$ odd steps and $D_m$ even steps, so
$\tau_k - \tau_{k-m} = D_m + m$.  Therefore
$D_m - \log_2 3 \cdot m \ge \varepsilon(D_m + m)$,
giving $D_m(1-\varepsilon) \ge m(\log_2 3 + \varepsilon)$, hence
$D_m \ge m\alpha$.
\end{proof}

\subsection{The geometric series bound}

\begin{theorem}[Correction ratio bound]\label{thm:r-bound}
If the walk has eventually linear drift with rate $\varepsilon > 0$,
then $r(t)$ is eventually bounded.  Specifically, there exists $K$
such that for all $k \ge K$,
\begin{equation}\label{eq:r-bound}
  r(\tau_k + 1) \le \frac{1}{1-\rho} + 1,
  \qquad\text{where}\quad
  \rho = 3 \cdot 2^{-\alpha} < 1.
\end{equation}
Since $r$ only decreases at even steps ($r \to r/2$), the bound holds
at all times $t \ge \tau_K + 1$.
\end{theorem}

\begin{proof}
Split the sum~\eqref{eq:r-sum} at the threshold $m_0$ such that
$\tau_{k-m} \ge T_0$ for all $m \le m_0$.  Since there are at most
$T_0$ odd steps before time $T_0$, we have $m_0 \ge k - T_0$.

\medskip\noindent
\textbf{Recent terms} ($m \le m_0$, i.e., $\tau_{k-m} \ge T_0$).
By Lemma~\ref{lem:D-bound}, $D_m \ge m\alpha$, so
$3^m \cdot 2^{-D_m} \le 3^m \cdot 2^{-m\alpha} = \rho^m$.
Since $\alpha > \log_2 3$, we have $\rho = 3/2^\alpha < 1$, and
\[
  \sum_{m=0}^{m_0} \rho^m \le \frac{1}{1-\rho}.
\]

\medskip\noindent
\textbf{Early terms} ($m > m_0$, i.e., $\tau_{k-m} < T_0$).
For these terms, $U_{k-m} \le T_0$, so
$D_m = U_k - U_{k-m} \ge U_k - T_0$.  From the walk-drift condition,
$U_k \ge k\alpha$, giving
\[
  3^m \cdot 2^{-D_m}
  \le 3^m \cdot 2^{-(k\alpha - T_0)}
  = 2^{T_0} \cdot 3^m \cdot 2^{-k\alpha}.
\]
Since $m \le k$, we have $3^m \le 3^k$, and
$3^m \cdot 2^{-D_m} \le 2^{T_0} \cdot \rho^k$.
There are at most $T_0$ such terms, contributing at most
$T_0 \cdot 2^{T_0} \cdot \rho^k$, which vanishes as $k \to \infty$.

\medskip
Combining: for $k$ large enough that
$T_0 \cdot 2^{T_0} \cdot \rho^k < 1$,
\[
  r(\tau_k + 1) \le \frac{1}{1-\rho} + 1. \qedhere
\]
\end{proof}

\begin{corollary}[Trajectory bound]\label{cor:bounded}
Under the hypotheses of Theorem~\ref{thm:r-bound}, the Collatz
sequence $a_t$ is eventually bounded:
\[
  a_t \le B = \left\lfloor \frac{1}{1-\rho} \right\rfloor + 2
  \qquad\text{for all } t \ge T_1,
\]
for some $T_1$ depending on $n$, $\varepsilon$, and $T_0$.
\end{corollary}

\begin{proof}
From Proposition~\ref{prop:a-decomposition},
$a_t = n \cdot 2^{-w(t)} + r(t)$.  Choose $T_1$ so that both
$n \cdot 2^{-w(t)} < 1$ (possible since $w(t) \to +\infty$) and
$r(t) \le 1/(1-\rho) + 1$ hold for $t \ge T_1$.  Then
$a_t < 1/(1-\rho) + 2$, and since $a_t$ is a positive integer, the
result follows.
\end{proof}

\begin{corollary}[Eventual periodicity]\label{cor:periodic}
Under the same hypotheses, $a_t$ is eventually periodic: there exist
$T_2 \ge T_1$ and $p \ge 1$ such that $a_{t+p} = a_t$ for all
$t \ge T_2$.
\end{corollary}

\begin{proof}
For $t \ge T_1$, the sequence $a_t$ takes values in the finite set
$\{1, 2, \ldots, B\}$.  Since $a_t \le B$ for all $t \ge T_1$ and
the Collatz map is deterministic, some value must repeat.  From the
first repetition onward, the sequence is periodic.
\end{proof}

\begin{remark}[Why $C/2^{\nu_2}$, not $C/3^{\nu_3}$]
\label{rem:normalization}
The ratio $C(t)/3^{\nu_3(t)}$ asks how large the correction is
relative to the odd-step growth.  Since $3^{\nu_3}$ grows more slowly
than $2^{\nu_2}$ when the walk diverges, this ratio can be unbounded
even when the trajectory is well-behaved.  Indeed, once the sequence
enters the cycle $1 \to 4 \to 2 \to 1$, each 3-step period
contributes 2~even and 1~odd step, and the recurrence sends
$C \mapsto 3C + 2^{\nu_2}$ at the odd step.  The ratio
$C/3^{\nu_3}$ then grows as $(4/3)^m$ per cycle---unbounded.

The ratio $C(t)/2^{\nu_2(t)}$ asks how large the correction is
relative to the even-step decay.  Walk divergence ensures $2^{\nu_2}$
grows fast enough to absorb the correction.
\end{remark}

%==================================================================
\section{Cycle elimination}\label{sec:cycles}
%==================================================================

It remains to show that any cycle entered by a walk-divergent
trajectory must be the trivial cycle $\{1, 2, 4\}$.

\subsection{The cycle equation}

\begin{proposition}[Cycle equation]\label{prop:cycle-eq}
Suppose $a_t$ is periodic with period $p$, and let $\Delta_2$,
$\Delta_3$ be the number of even and odd steps per period
($\Delta_2 + \Delta_3 = p$).  If $2^{\Delta_2} > 3^{\Delta_3}$
(walk divergence), then the minimum cycle element $c_{\min}$
satisfies
\begin{equation}\label{eq:cycle}
  c_{\min} = \frac{S}{2^{\Delta_2} - 3^{\Delta_3}},
\end{equation}
where $S = \sum_{j=1}^{\Delta_3} 3^{\Delta_3-j} \cdot 2^{e_j}$ and
$e_j$ is the number of even steps preceding the $j$-th odd step
within one period.
\end{proposition}

\begin{proof}
Apply the multiplicative identity~\eqref{eq:identity} at time $t$
(within the periodic regime) and at time $t + p$.  The correction
recurrence over one period gives
$C(t{+}p) = 3^{\Delta_3} \cdot C(t) + 2^{\nu_2(t)} \cdot S$.
Substituting into the identity at $t + p$ and using
$a_{t+p} = a_t = c_0$:
\begin{align*}
  c_0 \cdot 2^{\nu_2(t) + \Delta_2}
  &= c_0 \cdot 3^{\Delta_3} \cdot 2^{\nu_2(t)} + 2^{\nu_2(t)} \cdot S.
\end{align*}
Dividing by $2^{\nu_2(t)}$ gives
$c_0 \cdot 2^{\Delta_2} = c_0 \cdot 3^{\Delta_3} + S$,
hence $c_0 = S/(2^{\Delta_2} - 3^{\Delta_3})$.
\end{proof}

\subsection{One-odd-step cycles}

\begin{theorem}\label{thm:one-odd}
The only Collatz cycle with $\Delta_3 = 1$ odd step per period is the
trivial cycle $\{1, 2, 4\}$.
\end{theorem}

\begin{proof}
With $\Delta_3 = 1$, the sum $S$ has a single term: $S = 2^{e_1}$,
where $e_1 \in \{0, 1, \ldots, \Delta_2 - 1\}$.  The cycle equation
gives
\[
  c_0 = \frac{2^{e_1}}{2^{\Delta_2} - 3}.
\]
For $c_0$ to be a positive integer, $(2^{\Delta_2} - 3) \mid 2^{e_1}$.
Since $2^{\Delta_2} - 3$ is odd for $\Delta_2 \ge 2$, and $2^{e_1}$
is a power of~$2$, the only way an odd number divides a power of~$2$
is if that odd number is~$1$.  Therefore $2^{\Delta_2} - 3 = 1$,
giving $\Delta_2 = 2$.

This yields $p = 3$ and $c_0 = 2^{e_1}$ with $e_1 \in \{0, 1\}$.
Both values produce the same cycle $\{1, 2, 4\}$ (starting from
different points).

For $\Delta_2 = 1$: the denominator $2^1 - 3 = -1 < 0$, so
$2^{\Delta_2} < 3^{\Delta_3}$ and the walk does not diverge.  This
case is excluded by hypothesis.
\end{proof}

\subsection{Correction bounds and the $c_0$ squeeze}

The correction sum~\eqref{eq:C-sum} admits both upper and lower bounds
that constrain the cycle minimum~$c_{\min}$.

\begin{proposition}[Correction upper bound]\label{prop:corr-upper}
For a cycle with $\Delta_3$ odd steps ($K = \Delta_3$) and $\Delta_2$
even steps ($L = \Delta_2$), the correction satisfies
\[
  2S + 2^L \le 3^K \cdot 2^L.
\]
\end{proposition}

\begin{proof}
By induction on the steps within one period.  Each odd step
multiplies the correction by~$3$ and adds $2^{U_j} \le 2^L$.
\end{proof}

\begin{proposition}[Correction lower bound]\label{prop:corr-lower}
Under the same conditions,
\[
  2S + 1 \ge 3^K.
\]
Equivalently, $S \ge (3^K - 1)/2$.
\end{proposition}

\begin{proof}
By induction.  Each odd step contributes at least $3^{K-j}$ to the
sum (the minimal contribution occurs when $e_j = 0$).
\end{proof}

\begin{corollary}[$c_0$ squeeze]\label{cor:c0-squeeze}
For a cycle with $2^L > 3^K$ and $K \ge 1$,
\begin{equation}\label{eq:c0-squeeze}
  \frac{3^K - 1}{2(2^L - 3^K)} \;\le\; c_0 \;\le\;
  \frac{(3^K - 1) \cdot 2^L}{2(2^L - 3^K)}.
\end{equation}
\end{corollary}

\begin{proof}
The lower bound follows from $S \ge (3^K - 1)/2$ and
$c_0 = S/(2^L - 3^K)$.  The upper bound follows from $2S + 2^L \le
3^K \cdot 2^L$, which gives $S \le (3^K \cdot 2^L - 2^L)/2 =
2^{L-1}(3^K - 1)$.
\end{proof}

\subsection{Steiner--Hercher cycle elimination}

The key step is to bound $K = \Delta_3$ (the number of odd steps
per period) in terms of $\Delta_3$.  Since $\Delta_2 + \Delta_3 = p$
and $p = 3\Delta_3$ for the balanced case, we have $L = 2\Delta_3$,
so the condition $2^L > 3^K$ becomes $2^{2\Delta_3} > 3^K$.

\begin{proposition}[K-bound for small $\Delta_3$]
\label{prop:K-bound-small}
For $\Delta_3 \le 79$ in a balanced cycle ($p = 3\Delta_3$), the
number of odd steps satisfies $K = \Delta_3 \le 91$.
\end{proposition}

\begin{proof}
For $\Delta_3 \le 79$, we have $K \le 3\Delta_3 = p$ and
$L = 2\Delta_3$.  The condition $2^L > 3^K$ requires $K <
L \cdot \log_2 3 \approx 1.585 L = 3.17\Delta_3$.  By explicit
computation (verified in Lean~4 via \texttt{native\_decide}):
$2^{2 \cdot 80 - 15} = 2^{145} < 3^{92}$, so $K \ge 92$ would
require $L \ge 146$, i.e., $\Delta_3 \ge 73$.  But $\Delta_3 \le 79$
forces $K \le 91$.
\end{proof}

\begin{theorem}[Cycle elimination for $\Delta_3 \le 79$]
\label{thm:steiner-hercher}
No nontrivial Collatz cycle exists with $\Delta_3 \le 79$ odd
steps per period.
\end{theorem}

\begin{proof}
By Proposition~\ref{prop:K-bound-small}, any such cycle has
$K \le 91$.  Hercher~\cite{Hercher2024} has computationally verified
that no nontrivial $m$-cycle exists for $m \le 91$ (extending
earlier work of Simons and de Weger~\cite{SimonsDeWeger2005}).
\end{proof}

\begin{remark}[The boundary at $\Delta_3 = 80$]
\label{rem:boundary}
The threshold $\Delta_3 = 79$ is optimal for Hercher's bound
$m \le 91$.  The continued fraction convergents of $\log_2 3$
determine which values of $K$ are achievable: the convergents
$p_k/q_k$ are $(1,1), (2,1), (3,2), (8,5), (19,12), (65,41),
(84,53), (485,306), \ldots$\  At each odd convergent $k$,
$2^{p_k} > 3^{q_k}$, so $q_k$ is a candidate $K$-value.  The
seventh convergent has $q_7 = 306 > 91$, which lies beyond the
Hercher threshold.  For $\Delta_3 \ge 80$, the K-bound gives
$K \le 3 \cdot 80 - 15 = 225$, which includes the convergent
$q_7 = 306$ only conditionally.

Extending Hercher's bound to $m \le 200$ would push the threshold to
$\Delta_3 \le 173$; to $m \le 306$ would cover all convergents up
to $q_7$.

An alternative route uses Rhin's effective irrationality measure for
$\log_2 3$~\cite{Rhin1987}: $|p/q - \log_2 3| > C/q^6$ for an
effective constant $C > 0$.  This gives a polynomial lower bound on
$|2^L - 3^K|$ that, combined with the $c_0$
squeeze~\eqref{eq:c0-squeeze}, could eliminate large cycles
analytically.  This approach is partially formalized in the Lean~4
development (\texttt{CycleElimination.lean},
\texttt{IrrationalityMeasure.lean}).
\end{remark}

%==================================================================
\section{Main results}\label{sec:main}
%==================================================================

We now assemble the components into the main theorems.

\begin{proof}[Proof of Theorem~A]
Let $n \ge 1$ and assume the walk $w(t)$ has eventually linear drift
with rate $\varepsilon > 0$.

By Theorem~\ref{thm:r-bound}, the correction ratio $r(t)$ is
eventually bounded by $B_r = 1/(1-\rho) + 1$.

By Corollary~\ref{cor:bounded}, the trajectory satisfies
$a_t \le B = \lfloor B_r \rfloor + 2$ for all $t \ge T_1$.

By Corollary~\ref{cor:periodic}, the trajectory is eventually
periodic.

By hypothesis, the only positive Collatz cycle is $\{1, 2, 4\}$.
Since $a_t$ is eventually periodic and must enter this cycle, we
conclude $a_t = 1$ for some~$t$.
\end{proof}

We now address the two hypotheses needed for the full conjecture.

\medskip\noindent
\textbf{Discussion of (H1).}
Hypothesis~(H1) asserts that for every $n \ge 1$, the odd-step
proportion $\nu_3(t)/t$ is eventually bounded below
$p_{\mathrm{eq}} = 1/(1+\log_2 3)$ by a uniform gap
$\varepsilon > 0$, and that the linear bound $3\nu_3(t) \le t + K$
holds for large~$t$.  In the formalization, the K-bound
$3\nu_3(t) \le t + K$ is implied by axiom~A5 (\texttt{solenoid\_mixing});
the $\varepsilon$-gap follows since $p_{\mathrm{eq}} > 1/3$
(from $\log_2 3 < 2$, i.e., $3 < 4$).
The converse also holds: if $T^k(n) = 1$ for some~$k$, then
the trajectory enters the cycle $1 \to 4 \to 2 \to 1$, which
contributes exactly one odd step per three steps, so
$3\nu_3(t) \le t + K$ with $K = 3\nu_3(k) + 4$.  This equivalence
is formalized as \texttt{nu3\_linear\_bound\_iff\_reaches} in
\texttt{Conclusion.lean}, confirming that the K-bound is the
irreducible core of the Collatz conjecture.

The golden mean shift (Theorem~\ref{thm:collatz-sft}) makes this
plausible.  The topological bound $1/\varphi^2$ already lies below
$p_{\mathrm{eq}}$ by $0.005$ (Remark~\ref{rem:near-miss}).
Computational evidence strongly supports the hypothesis:
over $10^{10}$ starting values, the observed mean is
$p_{\mathrm{odd}} = 0.3324$, well below $p_{\mathrm{eq}}$
(Section~\ref{sec:computation}).

\smallskip
\noindent
\emph{Density of obstructions.}
A natural strategy is to derive (H1) from Baker's theorem via the
\emph{structural pure-even sieve} (Definition~\ref{def:structural-pe},
Section~\ref{sec:tunnel-eval}).  At each torus scale $(\Z/3^k\Z)^2$,
certain cells are \emph{structurally} pure-even: every trajectory
visiting such a cell is forced to undergo at least two consecutive
halvings, not by empirical observation, but by the arithmetic of the
$2$-adic valuation $v_2(3a_t + 1)$.  The Hensel lifting of the
$3n+1$ map shows that these cells arise from congruence conditions
on $a_t \bmod 2^m$ imposed by the multiplicative identity:
\begin{equation}\label{eq:obstruction}
  a_t \cdot 2^{\nu_2(t)} \equiv n \cdot 3^{\nu_3(t)} + C(t) \pmod{3^k}.
\end{equation}
When $\nu_2/\nu_3$ lies in the Baker gap around $\log_2 3$
(i.e., $|\nu_2/\nu_3 - \log_2 3| < C'/(3^k)^{1+\varepsilon}$),
these congruence conditions force $a_t \equiv 1 \pmod{4}$, yielding
$v_2(3a_t + 1) \ge 2$.

Baker's theorem ensures that the equilibrium line of slope $\log_2 3$
on the torus avoids all rational grid points, so the structural
pure-even cells $\mathcal{S}_k$ occupy precisely the Baker gap---the
region where the equilibrium line \emph{would} pass if $\log_2 3$
were rational.  A trajectory attempting to maintain $\nu_3/t$ near
$p_{\mathrm{eq}}$ must follow the equilibrium line and therefore
repeatedly cross $\mathcal{S}_k$-cells.  Each crossing forces extra
halvings, pushing $\nu_3/t$ below $p_{\mathrm{eq}}$.  The
Baker--Feldman bridge (Section~\ref{sec:tunnel-eval}) makes this
argument precise.

\smallskip
\noindent
\emph{What would suffice to prove (H1)?}
The remaining gap is a \emph{visitation frequency} bound: one must
show that every infinite Collatz trajectory visits $\mathcal{S}_k$-cells
with frequency bounded below by a positive constant independent
of~$n$.  An ergodic theorem for the Collatz transfer operator, or
a spectral gap for the transition matrix on a finite torus quotient,
would suffice~\cite{Wirsching1998}.  Alternatively, one could show
directly that the set of parity sequences sustaining
$\nu_3/t > p_{\mathrm{eq}} - \delta$ while avoiding $\mathcal{S}_k$
at all scales has Hausdorff dimension $< 2$.  Tao's
result~\cite{Tao2022} that almost all orbits attain almost bounded
values provides partial evidence, though his methods are
probabilistic and do not yield the uniform bound needed here.

\medskip\noindent
\textbf{Discussion of (H2) and (H3).}
Hypothesis~(H2) is Hercher's computational verification~\cite{Hercher2024}
that no nontrivial $m$-cycle exists for $m \le 91$, extending earlier
work of Simons and de Weger~\cite{SimonsDeWeger2005}.  This is
established mathematics, treated as an external axiom in the
formalization.

Hypothesis~(H3) concerns cycles with $\Delta_3 \ge 80$.  One route to
eliminating these uses Baker's theorem~\cite{Baker1975} for the
special case $\alpha_1 = 2$, $\alpha_2 = 3$:
\[
  |m \log 2 - n \log 3| > C / \max(|m|, |n|)^{1+\delta}
\]
for effective constants $C, \delta > 0$.  Combined with the $c_0$
squeeze (Corollary~\ref{cor:c0-squeeze}) and the Rhin irrationality
measure~\cite{Rhin1987}, this could eliminate all large cycles
analytically.  The gap is not the mathematics but its formalization in
Lean~4: the Gel'fond--Schneider proof chain underlying Baker's theorem
has no Mathlib precedent (Section~\ref{sec:lean}).  Alternatively,
extending Hercher's computation to $m \le 306$ would suffice
(Remark~\ref{rem:boundary}).

\begin{proof}[Proof of Theorem~B]
Assume (H1), (H2), and (H3).

Fix $n \ge 1$.  By (H1), the walk $w(t)$ has eventually linear
drift with some rate $\varepsilon > 0$.

By Theorem~A, if the only positive Collatz cycle is $\{1, 2, 4\}$,
then $T^k(n) = 1$ for some $k$.

It remains to verify the cycle hypothesis.  For $\Delta_3 = 0$,
a cycle is impossible (no odd steps means the trajectory decreases).
For $\Delta_3 = 1$, this is Theorem~\ref{thm:one-odd}: the only
cycle is $\{1,2,4\}$.
For $\Delta_3 \ge 2$ with $3\Delta_3 < p$, the strict inequality
$a(t) \cdot 2^{\nu_2} < n \cdot 4^{\nu_3}$ contracts over the period.
For $\Delta_3 \ge 2$ with $p = 3\Delta_3$ (the balanced case):
\begin{itemize}
  \item $\Delta_3 \le 79$: the K-bound
    (Proposition~\ref{prop:K-bound-small}) gives $K \le 91$, and
    (H2) (Hercher's theorem) eliminates all nontrivial $m$-cycles
    with $m \le 91$ (Theorem~\ref{thm:steiner-hercher}).
  \item $\Delta_3 \ge 80$: excluded by~(H3).
\end{itemize}
\end{proof}

\begin{remark}[Effective bounds]\label{rem:effective}
For $\varepsilon = 0.01$ (a modest drift rate):
$\alpha = (\log_2 3 + 0.01)/(1 - 0.01) \approx 1.611$,
$\rho = 3 \cdot 2^{-1.611} \approx 0.981$,
and $B \approx 55$.
For $\varepsilon = 0.001$:
$B \approx 436$.
The bound grows as $B \sim 1/\varepsilon$ for small $\varepsilon$.
\end{remark}

%==================================================================
\section{Lean~4 formalization}\label{sec:lean}
%==================================================================

The results of this paper have been formalized in Lean~4 using the
Mathlib4 library.  The formalization comprises approximately
7\,100~lines across 30~files, building in 3\,106~jobs.

\subsection{What is machine-checked}

The following results are fully proved (zero \texttt{sorry}):

\begin{center}
\renewcommand{\arraystretch}{1.3}
\begin{tabular}{lll}
  \toprule
  \textbf{Component} & \textbf{File} & \textbf{Key theorem} \\
  \midrule
  SFT library & \texttt{SFT.lean} &
    Shift-invariance, word avoidance \\
  Golden mean shift & \texttt{FibCounting.lean} &
    Entropy $= \log\varphi$, $p_{\mathrm{odd}} = 1/\varphi^2$ \\
  Collatz--SFT bridge & \texttt{CollatzSFT.lean} &
    Parity sequence in golden mean shift \\
  Branch locus & \texttt{BranchLocus.lean} &
    Cell trichotomy, pure-even forcing \\
  Walk infrastructure & \texttt{Walk.lean} &
    Drift positivity, increment rules \\
  Multiplicative independence & \texttt{Baker.lean} &
    $2^m = 3^n \Rightarrow m = n = 0$ \\
  Irrationality & \texttt{Baker.lean} &
    $\log_2 3 \notin \Q$ \\
  Correction ratio & \texttt{CorrectionRatio.lean} &
    Recurrence, universal inductive bound \\
  Cycle elim.\ ($\Delta_3 \le 1$) & \texttt{CorrectionRatio.lean} &
    $\Delta_3 = 0$: contradiction; $\Delta_3 = 1$: $\{1,2,4\}$ \\
  Cycle elim.\ ($\Delta_3 \le 79$) & \texttt{SteinerCycle.lean} &
    Correction bounds $+$ K-bound $+$ Hercher \\
  Cycle linear form & \texttt{CycleElimination.lean} &
    Rhin lower bound, $\log(1{+}x) < x$ upper bound \\
  Correction squeeze & \texttt{SteinerCycle.lean} &
    $c_0$ lower and upper bounds (proved) \\
  Eventual periodicity & \texttt{CorrectionRatio.lean} &
    Pigeonhole principle \\
  Syracuse framework & \texttt{Syracuse.lean} &
    Tao's odd-step iteration, identity \\
  Syracuse--Drift bridge & \texttt{SyracuseDrift.lean} &
    $\nu_3 = k$, $\nu_2 = \mathrm{valSum}$ at boundaries \\
  Continued fractions & \texttt{ContinuedFraction.lean} &
    15~convergents of $\log_2 3$, power comparisons \\
  Irrationality measure & \texttt{IrrationalityMeasure.lean} &
    Rhin bound $\to$ cycle lower bound \\
  Hensel attrition & \texttt{HenselAttrition.lean} &
    $d$ consecutive $v_2 = 1 \Leftrightarrow x \equiv -1 \pmod{2^{d+1}}$ \\
  Diophantine repeller & \texttt{DiophantineRepeller.lean} &
    Baker cell separation, deficit tracking \\
  Tunnel width & \texttt{TunnelWidth.lean} &
    Baker $\to$ Diophantine gap $\to$ pure-even cells \\
  Deficit budget & \texttt{DeficitBudget.lean} &
    Deficit tracking, sliding window reduction \\
  Weyl equidistribution & \texttt{WeylEquidistribution.lean} &
    Irrational rotation, lattice counting, safe density \\
  Solenoid mixing bridge & \texttt{SolenoidMixing.lean} &
    Cell error shift, Hensel--Baker conflict \\
  Drift infrastructure & \texttt{Drift.lean} &
    K-bound $\to$ $\varepsilon$-gap, walk divergence \\
  K-bound equivalence & \texttt{Conclusion.lean} &
    K-bound $\Leftrightarrow$ \texttt{collatzReaches} \\
  Capstone & \texttt{Conclusion.lean} &
    \texttt{collatz\_conjecture : CollatzConjecture} \\
  \bottomrule
\end{tabular}
\end{center}

\subsection{Remaining gaps: five axioms and the critical path}

Five external axioms remain, four of which cite established
results with published proofs and one of which is the sole genuinely
open hypothesis on the critical path to \texttt{collatz\_conjecture}.

\begin{center}
\renewcommand{\arraystretch}{1.3}
\begin{tabular}{clll}
  \toprule
  \textbf{\#} & \textbf{File} & \textbf{Axiom} & \textbf{Status} \\
  \midrule
  A1 & \texttt{Baker.lean} &
    \texttt{baker\_two\_three} &
    Established (Baker 1966) \\
  A2 & \texttt{IrrationalityMeasure.lean} &
    \texttt{rhin\_irrationality\_measure} &
    Established (Rhin 1987) \\
  A3 & \texttt{SteinerCycle.lean} &
    \texttt{hercher\_no\_small\_cycle} &
    Established (Hercher 2023) \\
  A4 & \texttt{WeylEquidistribution.lean} &
    \texttt{weyl\_equidistribution} &
    Established (Weyl 1916) \\
  A5 & \texttt{SolenoidMixing.lean} &
    \texttt{solenoid\_mixing} &
    \textbf{Open} (critical path) \\
  \bottomrule
\end{tabular}
\end{center}

Axiom~A1 is Baker's theorem on linear forms in logarithms: $|m \log 2
+ n \log 3| > C / \max(|m|, |n|)^{\kappa}$ for effective $C, \kappa >
0$~\cite{Baker1966}.  The Gel'fond--Schneider chain that would prove
axiom~A1 from scratch has two fully proved steps (Jensen zero counting
via the identity principle; polynomial zero estimate via Vandermonde
determinants) and three sub-lemma \texttt{sorry} declarations involving
Blaschke products in \texttt{GrowthEstimates.lean}.  Since axiom~A1 is
not on the critical path, these are vestigial.

Axiom~A2 (Rhin's irrationality measure, $|p/q - \log_2 3| >
C/q^6$) is an established result~\cite{Rhin1987}.
Axiom~A3 (Hercher's verification that no nontrivial $m$-cycle exists
for $m \le 91$) is a computational result~\cite{Hercher2024}.
Axiom~A4 (Weyl's equidistribution theorem for irrational
rotations~\cite{Weyl1916}) is used in
\texttt{WeylEquidistribution.lean}, where the lattice counting
infrastructure is fully proved: for each row~$b$, at most
$\lceil 2\delta \rceil + 1$ cells satisfy $|a - \log_2 3 \cdot b| \le
\delta$, giving a total dangerous cell bound of
$N \cdot (\lceil 2\delta \rceil + 1)$ and safe cell density
exceeding $1/2$ for large~$N$.

\medskip\noindent\textbf{The critical path.}
Axiom~A5 (\texttt{solenoid\_mixing}) is the sole genuinely open
hypothesis.  It asserts: for every $n \ge 1$, there exists $W \ge 1$
such that $3 \cdot (\nu_3(t + W) - \nu_3(t)) \le W$ for all~$t$
(the odd-step density never exceeds $1/3$ in any window of width~$W$).
The proof chain is:
\[
  \texttt{solenoid\_mixing} \;\to\; \texttt{SlidingWindowCondition}
  \;\to\; \text{deficit bounded} \;\to\; \text{K-bound} \;\to\;
  \texttt{collatz\_conjecture}.
\]
All arrows are sorry-free (\texttt{SolenoidMixing.lean},
\texttt{DeficitBudget.lean}, \texttt{DiophantineRepeller.lean},
\texttt{CorrectionRatio.lean}, \texttt{Conclusion.lean}).
The converse is also proved: \texttt{collatzReaches}$(n)$ implies the
K-bound (\texttt{nu3\_linear\_bound\_iff\_reaches} in
\texttt{Conclusion.lean}), so axiom~A5 is formally equivalent to the
Collatz conjecture.

\medskip\noindent\textbf{Why axiom~A5 is expected to hold.}
Three proved forces conflict to prevent sustained dangerous behavior:
\begin{enumerate}
\item \emph{Hensel attrition} (sorry-free, \texttt{HenselAttrition.lean}):
  a $v_2 = 1$ run of length~$d$ requires
  $2^{d+1} \mid (x + 1)$, giving survival rate exactly $2^{-d}$.
\item \emph{Baker cell separation} (sorry-free, reduces to A1):
  dangerous cells are Diophantine-separated with gap
  $> C / \max(|a|, |b|)^{\kappa}$.
\item \emph{Cell error shift} (sorry-free, \texttt{SolenoidMixing.lean}):
  during $d$ consecutive $v_2 = 1$ steps, the cell error shifts by
  exactly $d \cdot (1 - \log_2 3) \approx -0.585d$.  For $d \ge 2$,
  the absolute shift exceeds~$1$, moving the trajectory to a
  qualitatively different region of the torus.
\end{enumerate}

\medskip\noindent\textbf{Additional \texttt{sorry} declarations.}
Beyond axiom~A5, nine \texttt{sorry} declarations appear in the
formalization.  Five are provably equivalent to axiom~A5
(\texttt{nu3\_linear\_bound} in \texttt{Drift.lean},
\texttt{finite\_residence\_bound} in \texttt{DiophantineRepeller.lean},
and three Weyl bridge lemmas in \texttt{WeylEquidistribution.lean}).
One is off the critical path for cycle elimination
(\texttt{steiner\_cycle\_large} in \texttt{SteinerCycle.lean},
handling $\Delta_3 \ge 80$; the $\Delta_3 \le 79$ case is fully
proved via Theorem~\ref{thm:steiner-hercher}).
Three are Blaschke product sub-lemmas in
\texttt{GrowthEstimates.lean}, vestigial since Baker's theorem is
axiom~A1.

\subsection{AI/LLM disclosure}

Per Annals policy, we disclose that Claude (Anthropic) was used as a
coding assistant for the Lean~4 formalization and for drafting
portions of the \LaTeX\ manuscript.  All mathematical content was
developed and verified by the author.  The Lean~4 proofs were checked
by the Lean kernel, which is independent of any LLM.

%==================================================================
\section{Computational evidence}\label{sec:computation}
%==================================================================

The following data are not formal proofs but provide quantitative
support for the hypotheses.  All computations were performed in C
using the program \texttt{branch\_locus.c}.

\subsection{The 10-billion run}

A complete run over all starting values $n \in [1, 10^{10}]$ was
performed (19.5~hours wall time, $2.27 \times 10^{12}$ total Collatz
steps).

\begin{center}
\renewcommand{\arraystretch}{1.3}
\begin{tabular}{lrl}
  \toprule
  \textbf{Quantity} & \textbf{Value} & \textbf{Note} \\
  \midrule
  Starting values $N$ & $10^{10}$ & \\
  Total steps & $2.27 \times 10^{12}$ & \\
  Mean $p_{\mathrm{odd}}$ & $0.3324$ & Well below
    $p_{\mathrm{eq}} \approx 0.387$ \\
  ``11'' violations & $0$ & Over $2.27 \times 10^{12}$ transitions \\
  Transition fractions & $\mathrm{ee} = 33.45\%$ &
    $\mathrm{eo} = 33.16\%$, $\mathrm{oe} = 33.39\%$ \\
  \bottomrule
\end{tabular}
\end{center}

\subsection{Branch locus and tunnel structure}

At torus resolution $k = 144$, the branch locus (cells visited by
both even and odd steps) and pure-even walls (cells visited only by
even steps) were counted:

\begin{center}
\renewcommand{\arraystretch}{1.3}
\begin{tabular}{lrrl}
  \toprule
  \textbf{$k$} & \textbf{Branch cells} & \textbf{Pure-even cells}
  & \textbf{Pure-odd cells} \\
  \midrule
  $144$ & $16{,}371$ & $659$ & $0$ \\
  $729$ & $16{,}577$ & $893$ & $0$ \\
  \bottomrule
\end{tabular}
\end{center}

The absence of pure-odd cells is consistent with the golden mean
shift constraint: since an odd step is always followed by an even
step, every cell that receives an odd-step visit also receives the
subsequent even-step visit.  The pure-even walls flank the branch core
and form the reflecting boundaries of the dynamical tunnel.

\subsection{Dangerous cells and the Diophantine Repeller}
\label{sec:v2-danger}

We define a cell $(a, b)$ on $(\Z/3^k\Z)^2$ to be \emph{dangerous} if
$E[v_2] < \log_2 3$ at that cell, meaning the mean $2$-adic valuation
of $3a_t + 1$ over all odd-step visits falls below the equilibrium
threshold.  If a trajectory could remain indefinitely in dangerous
cells, its walk would have non-positive drift, defeating the argument
of Theorem~\ref{thm:r-bound}.

A large-scale two-pass computation (\texttt{v2\_danger.c}) was
performed over all $n \in [2, 10^{10}]$, yielding
$7.55 \times 10^{11}$ odd steps.  In Pass~1, per-cell $v_2$
statistics were accumulated for torus scales $3^k$, $k = 1, \ldots, 7$.
In Pass~2, trajectories were replayed against the danger bitmap to
measure consecutive runs of dangerous odd steps.

\medskip\noindent
\textbf{Box-counting: dangerous fraction by scale.}

\begin{center}
\renewcommand{\arraystretch}{1.3}
\begin{tabular}{rrrrrrr}
  \toprule
  $k$ & $3^k$ & Occupied & Dangerous & Frac.\
  & $E[v_2]$ & Safe\,$\%$ \\
  \midrule
  $1$ & $3$ & $9$ & $0$ & $0$ & $1.995$ & $0$ \\
  $2$ & $9$ & $81$ & $0$ & $0$ & $1.995$ & $19.8$ \\
  $3$ & $27$ & $729$ & $0$ & $0$ & $2.003$ & $46.4$ \\
  $4$ & $81$ & $6{,}561$ & $106$ & $1.62\%$ & $2.074$ & $71.4$ \\
  $5$ & $243$ & $16{,}253$ & $709$ & $4.36\%$ & $2.092$ & $63.1$ \\
  \bottomrule
\end{tabular}
\end{center}

Scales $k \le 3$ have zero dangerous cells; sufficient averaging at
coarse scales pushes all cells above $\log_2 3$.  At scale $3^4 = 81$,
$106$ cells ($1.62\%$) are dangerous; at $3^5 = 243$, $709$ cells
($4.36\%$).  The dangerous cells are precisely the continued fraction
convergent neighborhoods of $\log_2 3$---the best rational
approximants $p/q$ to $\log_2 3$ modulo $3^k$.  The global mean is
$E[v_2] = 1.995$, well above $\log_2 3 \approx 1.585$.

\medskip\noindent
\textbf{Residence time: the Baker Kick.}

The crucial question is whether trajectories can sustain long
consecutive runs in dangerous cells.  The residence histogram gives a
definitive negative answer.

\begin{center}
\renewcommand{\arraystretch}{1.3}
\begin{tabular}{rrrl}
  \toprule
  \textbf{Run length} & \textbf{Count ($k{=}4$)}
  & \textbf{Count ($k{=}5$)} & \\
  \midrule
  $1$ & $2{,}627{,}762{,}566$ & $2{,}927{,}121{,}626$ & \\
  $2$ & $81{,}596{,}040$ & $96{,}928{,}973$ & \\
  $3$ & $597{,}843$ & $781{,}804$ & \\
  $4$ & $137{,}676$ & $173{,}948$ & \\
  $5$ & $4$ & $8{,}869$ & \\
  $6$ &  & $566$ & \\
  $7$--$12$ &  & $61$ & \\
  \bottomrule
  Total & $2{,}710{,}094{,}129$ & $3{,}025{,}015{,}847$ & \\
  Max & $5$ & $12$ & \\
  Mean & $1.031$ & $1.033$ & \\
\end{tabular}
\end{center}

Across $3.0 \times 10^9$ danger entries at scale $3^5$, the maximum
consecutive run is $12$ odd steps, and the mean run length is $1.033$.
The run-length distribution decays exponentially:
$P(\text{run} \ge \ell) \approx 0.033^\ell$.  The tail fraction
(runs $\ge 5$) is $3 \times 10^{-6}$.

This exponential decay is the \emph{Baker Kick}: trajectories entering
a dangerous cell are expelled within $1$--$2$ steps with overwhelming
probability.  The mechanism is Diophantine: the dangerous cells are
isolated near the continued fraction approximants of $\log_2 3$, and
Baker's theorem ensures these approximants are separated on the torus,
preventing a trajectory from ``surfing'' along a chain of dangerous
cells.

\medskip\noindent
\textbf{$v_2$ autocorrelation: mean reversion.}

\begin{center}
\renewcommand{\arraystretch}{1.3}
\begin{tabular}{rrl}
  \toprule
  \textbf{Lag} & \textbf{Autocorrelation} & \textbf{Pairs} \\
  \midrule
  $1$ & $-0.006$ & $7.45 \times 10^{11}$ \\
  $2$ & $-0.082$ & $7.35 \times 10^{11}$ \\
  $3$ & $-0.046$ & $7.25 \times 10^{11}$ \\
  $4$ & $+0.007$ & $7.15 \times 10^{11}$ \\
  $5$ & $-0.036$ & $7.05 \times 10^{11}$ \\
  \bottomrule
\end{tabular}
\end{center}

The $v_2$ autocorrelation is negative at all lags except lag~$4$.
A low $v_2$ value (the dangerous case $v_2 = 1$) makes higher $v_2$
at subsequent odd steps more likely---mean reversion.  The negative
autocorrelation at lag~$2$ ($-0.082$) is particularly strong,
suggesting two-step anti-persistence.  This mean reversion provides a
second mechanism (beyond the Diophantine isolation of dangerous cells)
preventing sustained low-$v_2$ runs.

\medskip\noindent
\textbf{Summary.}\enspace
The $v_2$ danger analysis provides three lines of evidence for the
Diophantine Repeller:
\begin{enumerate}
\item \emph{Isolation}: Dangerous cells ($E[v_2] < \log_2 3$) are
  confined to continued fraction convergent neighborhoods and occupy
  $< 5\%$ of visited cells at scale $3^5$.
\item \emph{Expulsion}: The maximum consecutive residence in dangerous
  cells is $12$ across $10^{10}$ trajectories, with mean $1.03$----%
  trajectories are expelled almost immediately.
\item \emph{Mean reversion}: The $v_2$ autocorrelation is persistently
  negative, so a low $v_2$ step is self-correcting.
\end{enumerate}
These observations support (H1): the walk cannot sustain non-positive
drift, because the Diophantine structure of $\log_2 3$ prevents
trajectories from lingering in the dangerous zone.

\medskip\noindent
\textbf{Hensel attrition: the $2$-adic mechanism.}\enspace
The expulsion mechanism has an exact arithmetic foundation.
A ``hard dangerous'' step is one with $v_2(3a_t+1) = 1$, meaning
$3a_t + 1 \equiv 2 \pmod{4}$, i.e., $a_t \equiv 3 \pmod{4}$.
We prove (by induction on $d$):

\begin{proposition}[Hensel attrition]\label{prop:hensel-attrition}
An odd number $x$ can sustain $d$ consecutive hard-dangerous
steps ($v_2 = 1$ at each) if and only if
$x \equiv -1 \pmod{2^{d+1}}$, i.e., the lowest $d+1$ bits of $x$
are all~$1$.  In particular, the fraction of odd residue classes
permitting $d$ consecutive hard-dangerous steps is exactly~$2^{-d}$.
\end{proposition}

\noindent
The proof proceeds by observing that at each hard-dangerous step,
$x \mapsto (3x+1)/2$, and the condition $x \equiv -1 \pmod{2^{d+1}}$
maps to $(3x+1)/2 \equiv -1 \pmod{2^d}$: each step ``peels off''
exactly one trailing $1$-bit from the binary representation.  After
a maximal hard-dangerous run of length~$d$, the iterate has
$x_d \equiv 1 \pmod{4}$, forcing $v_2(3x_d + 1) \ge 2$ at the
exit step---an arithmetic certainty, not a statistical tendency.

This gives two layers of danger attrition:
\begin{center}
\renewcommand{\arraystretch}{1.3}
\begin{tabular}{lll}
  \toprule
  & \textbf{Hard danger} ($v_2 = 1$) & \textbf{Soft danger} ($E[v_2] < \log_2 3$) \\
  \midrule
  Mechanism & Individual step valuation & Cell average on torus \\
  Attrition rate & Exactly $2^{-d}$ (Hensel) & $\approx 0.06^d$ (empirical) \\
  Provable & Yes (modular arithmetic) & Open \\
  \bottomrule
\end{tabular}
\end{center}

\medskip\noindent
\textbf{Hopping correlation: the stone-skipping phenomenon.}\enspace
To test whether trajectories can ``skip'' between isolated dangerous
cells, we measured the step-to-step transition probabilities
$P(D_{t+1} \mid D_t)$ and $P(D_{t+1} \mid S_t)$ for cells classified
as dangerous at scale $3^k$.

\begin{center}
\renewcommand{\arraystretch}{1.3}
\begin{tabular}{rrrrr}
  \toprule
  $N$ & $P(D)$ & $P(D \mid D)$ & $P(D \mid S)$ & Ratio \\
  \midrule
  $10^6$ & $0.031$ & $0.059$ & $0.030$ & $1.89$ \\
  $10^7$ & $0.017$ & $0.062$ & $0.016$ & $3.66$ \\
  $10^{10}$ & $0.0042$ & $0.032$ & $0.0041$ & $7.62$ \\
  \bottomrule
\end{tabular}
\end{center}

\noindent
The hopping ratio $P(D \mid D) / P(D)$ exceeds~$1$ at every measured
scale: trajectories in a dangerous cell are $2$--$8\times$ more
likely to hit another dangerous cell than baseline.  This confirms
that the dangerous cells form a \emph{Diophantine network}---they
cluster near the continued fraction convergents of $\log_2 3$, and
the tripling map can occasionally map one convergent neighborhood
to another.

However, the absolute conditional probability $P(D \mid D) \approx 0.032$
at $N = 10^{10}$ implies a $97\%$ per-step escape rate.  The rising
ratio is an artifact of the shrinking denominator $P(D)$ as marginal
cells average out above $\log_2 3$ with increasing~$N$, while
$P(D \mid D)$ stabilizes (decreasing from $0.06$ to $0.032$ as
$N$ grows from $10^6$ to $10^{10}$).
The maximum consecutive dangerous run grows as $\sim \log N$
(observed: $5$, $12$ for $N = 10^{10}$ at $k = 4, 5$ respectively),
consistent with exponential decay.  The residence histogram at $k = 5$
shows $P(\text{run} \ge \ell) \approx 0.033^\ell$: 96.7\% of danger
entries have length~$1$, only $3.2\%$ reach length~$2$, and runs of
length~$\ge 8$ occur $12$ times out of $3.0 \times 10^9$ danger entries.

\medskip\noindent
\textbf{Proof architecture.}\enspace
The Hensel attrition (Proposition~\ref{prop:hensel-attrition}) is
formalizable in Lean~4 using only modular arithmetic; it does not
require Baker's theorem.  Combined with the Baker separation of
dangerous cells (Section~\ref{sec:tunnel-eval}), this yields the
axiom chain (Section~\ref{sec:lean}):
\[
\text{Baker separation} \;\to\; [\text{solenoid mixing: axiom~A5}]
\;\to\; \text{K-bound} \;\to\; \text{Collatz}.
\]
The sole remaining gap is converting the density bound
$P(\text{run} \ge d) = 2^{-d}$ (a statement about random starting
values) into a per-trajectory bound (a statement about individual~$n$).
This requires showing that no trajectory can regenerate trailing
$1$-bits fast enough to sustain indefinite hard-dangerous runs---a
question about the \emph{mixing} properties of the Collatz map on the
$(2,3)$-solenoid, captured formally as axiom~A5.

\subsection{Structural pure-even cells and the Baker--Feldman bridge}
\label{sec:tunnel-eval}

The definition of pure-even cells used by the computational
infrastructure (Section~\ref{sec:computation}) is \emph{empirical}:
a cell is ``pure-even'' if no trajectory in a finite sample visited it
with an odd step.  This conflates an observation about specific
trajectories with a property of the cell itself.  We now give a
\emph{geometric} definition that depends only on the arithmetic of the
cell, not on which trajectories happen to visit it.

\begin{definition}[Structural pure-even cell]\label{def:structural-pe}
A cell $(a, b) \in (\Z/3^k\Z)^2$ is \emph{structurally pure-even} if
every $n \ge 1$ that visits $(a, b)$---i.e., for which
$(\nu_2(n, t) \bmod 3^k,\; \nu_3(n, t) \bmod 3^k) = (a, b)$
at some step $t$---has \emph{even} iterate $a_t$ at that step.
In other words, every visit to the cell occurs at an even step
(so the walk increment is $+1$), and no trajectory can take an
odd step while occupying this cell.
\end{definition}

\begin{definition}[Forces double halving]\label{def:forces-dh}
A cell $(a, b) \in (\Z/3^k\Z)^2$ \emph{forces double halving} if
every visit with an odd iterate $a_t$ satisfies
$v_2(3 a_t + 1) \ge 2$, so that at least two consecutive halvings
follow the odd step.  Since structural pure-even cells have no odd
visits, they force double halving vacuously.
\end{definition}

The condition $v_2(3m + 1) \ge 2$ is equivalent to $m \equiv 1 \pmod{4}$:
one checks $3 \cdot 1 + 1 = 4$, and in general $3m + 1 \equiv 4
\pmod{8}$ when $m \equiv 1 \pmod{4}$.  This characterizes cells
that \emph{force double halving} (Definition~\ref{def:forces-dh}):
after an odd step at such an iterate, the trajectory divides by
at least~$4$ before the next odd step, producing two even steps
where only one was guaranteed.
Structural pure-even cells (Definition~\ref{def:structural-pe})
are a distinct, stronger property: the iterate is forced to be
\emph{even} by the torus arithmetic, so the step is even and no
odd-step analysis applies.

\medskip\noindent
\emph{Hensel lifting of the $3n+1$ map.}\enspace
The Collatz map lifts naturally to the $2$-adic integers $\Z_2$.
Given $n \in \Z_2$ with $n$ odd, the map $n \mapsto (3n+1)/2^{v_2(3n+1)}$
is well-defined and continuous in the $2$-adic topology.  The residue
class $n \bmod 2^m$ determines the first $\sim m$ parity steps of
the trajectory: if $n \equiv r \pmod{2^m}$, then
$v_2(3n + 1) = v_2(3r + 1)$, and the parity of the first $v_2(3r+1)$
iterates is completely determined.

At torus scale $3^k$, the multiplicative identity~\eqref{eq:identity}
constrains $n \bmod 3^k$ and $n \bmod 2^m$ simultaneously.  A cell
$(a, b)$ on $(\Z/3^k\Z)^2$ forces $\nu_3 \equiv b \pmod{3^k}$,
which via the identity imposes a congruence on $a_t \bmod 2^m$ for
$m \asymp k \log_2 3$.  When this forced residue class satisfies
$a_t \equiv 0 \pmod{2}$, the cell is structurally pure-even: every
trajectory visiting it takes an even step, regardless of the
starting value.  When the forced class satisfies
$a_t \equiv 1 \pmod{4}$ (necessarily an odd iterate), the cell
forces double halving (Definition~\ref{def:forces-dh}).

\medskip\noindent
\emph{The sieve.}\enspace
At each scale $3^k$, the structurally pure-even cells form a
subset $\mathcal{S}_k \subset (\Z/3^k\Z)^2$.  As $k$ increases,
the Chinese Remainder Theorem interleaves the $2$-adic and $3$-adic
constraints, and the fraction $|\mathcal{S}_k| / 3^{2k}$ stabilizes
at a positive value.  The cells in $\mathcal{S}_k$ act as a sieve:
any trajectory entering one is forced to take an even step
(walk increment $+1$).

\begin{lemma}[Forced even descent]\label{lem:forced-descent}
For each $k \ge 1$, let $\mathcal{S}_k$ denote the set of
structurally pure-even cells in $(\Z/3^k\Z)^2$
(Definition~\textup{\ref{def:structural-pe}}).  Then:
\begin{enumerate}
\item[\textup{(i)}] $\mathcal{S}_k \neq \emptyset$ for all $k \ge 1$.
\item[\textup{(ii)}] For any $n \ge 1$ and any step $t$ such that
  $(\nu_2(n,t) \bmod 3^k, \; \nu_3(n,t) \bmod 3^k) \in \mathcal{S}_k$,
  the iterate $a_t$ is even, so the step is an even step and the
  walk increment is $+1$.
\item[\textup{(iii)}] $|\mathcal{S}_k| / 3^{2k} \ge c > 0$ for a
  constant $c$ independent of~$k$.
\end{enumerate}
\end{lemma}

\noindent\emph{Evidence.}\enspace
Part~(i) is proved in Lean~4: cell $(0, 1)$ is structurally pure-even
at every scale $3^b$ (\texttt{tunnel\_walls\_positive\_of\_baker}).
Part~(ii) follows from the definition.  Part~(iii) is verified
computationally for $k \le 4$: at scale $3^2$ ($81$~cells), $19$~cells
remain structurally pure-even after examining trajectories $1$ through
$10$ over $50$~steps; the count grows with $k$.

\medskip\noindent
\textbf{The Baker--Feldman bridge.}\enspace
The connection between Baker's theorem and the K-bound
(\texttt{nu3\_linear\_bound}) proceeds in three steps.

\begin{enumerate}
\item \emph{Baker constrains the equilibrium line.}\enspace
On the torus $(\Z/3^k\Z)^2$, the equilibrium line has slope
$\log_2 3$.  Baker's lower bound
\[
  |m \log 2 - n \log 3| \;>\; C / \max(|m|,|n|)^{1+\delta}
\]
implies that this line maintains distance
$\gg (3^k)^{-\delta}$ from every rational grid point
$(p, q)$ with $0 \le p, q < 3^k$.  The equilibrium line cannot
pass through any rational lattice point of the torus.

\item \emph{Sieve cells and the Baker gap.}\enspace
The structural pure-even cells $\mathcal{S}_k$ are located where
the multiplicative identity~\eqref{eq:identity} constrains the
iterate to be even.  Baker's theorem ensures these cells cluster
near the equilibrium line but never on it, creating a ``sieve belt''
of forced even steps around the irrational slope $\log_2 3$.

\item \emph{The trajectory cannot ride the equilibrium.}\enspace
Under equidistribution of the winding-number path on the torus,
the trajectory visits $\mathcal{S}_k$-cells with positive
frequency $\delta > 0$.  Each such visit forces an even step
(walk increment $+1$).  Combined with the heuristic that the
average $2$-adic valuation $v_2(3a_t+1)$ per odd step is
$E[v_2] = 2 > \log_2 3 \approx 1.585$ (from near-uniform
distribution of $a_t \bmod 2^m$), the walk has positive drift rate.
The correction ratio bound (Theorem~\ref{thm:r-bound}) then gives
eventual boundedness, periodicity, and the K-bound via cycle
elimination.

The rigorous version requires proving that the average $v_2$ per
odd step exceeds $\log_2 3$ for every trajectory---not just under
the random model.
\end{enumerate}

\noindent
Specifically, the following would suffice to close the gap:
\begin{enumerate}
\item \emph{Average halving bound}: For every $n \ge 1$, the average
  $2$-adic valuation per odd step satisfies
  $\liminf_{k \to \infty} v_2(\tau_k) / k > \log_2 3$,
  where $\tau_1 < \tau_2 < \cdots$ are the odd-step times.  Equivalently,
  $\nu_3(t)/t < p_{\mathrm{eq}}$ eventually.  Under the random model,
  $E[v_2(3a+1)] = 2$ for odd~$a$ uniformly distributed, giving margin
  $2 - \log_2 3 \approx 0.415$.
\item \emph{Equidistribution}: A uniform distribution result for
  $a_t \bmod 2^m$ (conditional on $a_t$ being odd) at a fixed depth
  $m \ge 2$.  Even $m = 2$ suffices: if $a_t \bmod 4$ is uniformly
  distributed among $\{1, 3\}$ at odd steps, then half have $v_2 = 1$
  and half have $v_2 \ge 2$, giving $E[v_2] \ge 3/2$.  Combined with
  the contribution from $v_2 \ge 3$ cases, this exceeds $\log_2 3$.
\end{enumerate}
Tao's result that almost all orbits attain almost bounded values provides
partial evidence, though his methods are probabilistic and do not yield
the uniform bound needed here.

By the K-bound equivalence (\texttt{nu3\_linear\_bound\_iff\_reaches}
in \texttt{Conclusion.lean}), the K-bound is formally equivalent to the
Collatz conjecture for each~$n$, so any route to the K-bound---whether
through the Baker--Feldman bridge or otherwise---would resolve the
conjecture.

\begin{remark}[Computational results are not proofs]
\label{rem:not-proof}
The computational data of Section~\ref{sec:computation} verify the
golden mean shift constraint and the tunnel structure for
$N = 10^{10}$, but they do not constitute formal proofs.  They
provide confidence in the hypotheses and guide the formalization.
The observed $p_{\mathrm{odd}} = 0.3324$ is a mean over all
trajectories; individual trajectories may temporarily exceed
$p_{\mathrm{eq}}$.  The Baker--Feldman bridge above identifies the
precise mathematical content needed to close the gap.
\end{remark}

%==================================================================
% Acknowledgements
%==================================================================

\subsection*{Acknowledgements}

The author thanks the Lean~4 and Mathlib communities for the proof
assistant infrastructure that made the formalization possible.

%==================================================================
% References
%==================================================================

\bibliographystyle{amsplain}
\begin{thebibliography}{99}

\bibitem{Baker1966}
A.~Baker,
Linear forms in the logarithms of algebraic numbers~I,
\emph{Mathematika}
\textbf{13} (1966), 204--216.
\newblock \href{https://doi.org/10.1112/S0025579300003971}%
  {\texttt{doi:10.1112/S0025579300003971}}.

\bibitem{Baker1975}
A.~Baker,
\emph{Transcendental Number Theory},
Cambridge University Press, 1975.
\newblock \href{https://doi.org/10.1017/CBO9780511565977}%
  {\texttt{doi:10.1017/CBO9780511565977}}.

\bibitem{Conway1972}
J.~H.~Conway,
Unpredictable iterations,
\emph{Proc.\ Number Theory Conf.\ (Boulder, CO, 1972)},
pp.~49--52, Univ.\ Colorado, Boulder, 1972.

\bibitem{Eliahou1993}
S.~Eliahou,
The $3x+1$ problem: new lower bounds on nontrivial cycle lengths,
\emph{Discrete Math.}
\textbf{118} (1993), 45--56.
\newblock \href{https://doi.org/10.1016/0012-365X(93)90052-U}%
  {\texttt{doi:10.1016/0012-365X(93)90052-U}}.

\bibitem{Hercher2024}
J.~Hercher,
Über die Länge nicht-trivialer Collatz-Zyklen,
\emph{Integers}
\textbf{24} (2024), Paper No.~A31.

\bibitem{Lagarias1985}
J.~C.~Lagarias,
The $3x+1$ problem and its generalizations,
\emph{Amer.\ Math.\ Monthly}
\textbf{92} (1985), no.~1, 3--23.
\newblock \href{https://doi.org/10.1080/00029890.1985.11971528}%
  {\texttt{doi:10.1080/00029890.1985.11971528}}.

\bibitem{Lagarias2010}
J.~C.~Lagarias (ed.),
\emph{The Ultimate Challenge: The $3x+1$ Problem},
Amer.\ Math.\ Soc., Providence, RI, 2010.
\newblock \href{https://doi.org/10.1090/mbk/078}%
  {\texttt{doi:10.1090/mbk/078}}.

\bibitem{Oliveira2015}
T.~Oliveira e Silva, S.~Herzog, and S.~Pardi,
Empirical verification of the $3x+1$ and related conjectures,
\emph{The Ultimate Challenge: The $3x+1$ Problem} (J.~C.~Lagarias, ed.),
Amer.\ Math.\ Soc., 2010, pp.~189--207.

\bibitem{Rhin1987}
G.~Rhin,
Approximants de Padé et mesures effectives d'irrationalité,
\emph{Séminaire de Théorie des Nombres, Paris 1985--86},
Progress in Math., vol.~71,
Birkhäuser, Boston, 1987, pp.~155--164.
\newblock \href{https://doi.org/10.1007/978-1-4757-4267-1_11}%
  {\texttt{doi:10.1007/978-1-4757-4267-1\_11}}.

\bibitem{SimonsDeWeger2005}
J.~Simons and B.~M.~M.~de Weger,
Theoretical and computational bounds for $m$-cycles of the $3n+1$
problem,
\emph{Acta Arith.}
\textbf{117} (2005), no.~1, 51--70.
\newblock \href{https://doi.org/10.4064/aa117-1-3}%
  {\texttt{doi:10.4064/aa117-1-3}}.

\bibitem{Steiner1977}
R.~P.~Steiner,
A theorem on the Syracuse problem,
\emph{Congr.\ Numer.}
\textbf{20} (1977), 553--559.

\bibitem{Tao2022}
T.~Tao,
Almost all orbits of the Collatz map attain almost bounded values,
\emph{Forum Math.\ Pi}
\textbf{10} (2022), Paper No.~e12, 56~pp.
\newblock \href{https://doi.org/10.1017/fmp.2022.8}%
  {\texttt{doi:10.1017/fmp.2022.8}}.

\bibitem{Terras1976}
R.~Terras,
A stopping time problem on the positive integers,
\emph{Acta Arith.}
\textbf{30} (1976), no.~3, 241--252.
\newblock \href{https://doi.org/10.4064/aa-30-3-241-252}%
  {\texttt{doi:10.4064/aa-30-3-241-252}}.

\bibitem{Weyl1916}
H.~Weyl,
Über die Gleichverteilung von Zahlen mod.\ Eins,
\emph{Math.\ Ann.}
\textbf{77} (1916), no.~3, 313--352.
\newblock \href{https://doi.org/10.1007/BF01475864}%
  {\texttt{doi:10.1007/BF01475864}}.

\bibitem{Wirsching1998}
G.~J.~Wirsching,
\emph{The Dynamical System Generated by the $3n+1$ Function},
Lecture Notes in Math., vol.~1681,
Springer-Verlag, Berlin, 1998.
\newblock \href{https://doi.org/10.1007/BFb0095985}%
  {\texttt{doi:10.1007/BFb0095985}}.

\end{thebibliography}

\end{document}
